
\documentclass[10pt,a4paper]{article}

\usepackage[utf8]{inputenc}
\usepackage[T1]{fontenc}
\usepackage[french]{babel}
\usepackage{amsmath,amssymb}
\usepackage{geometry}
\usepackage{enumitem}
\usepackage{multicol}
\usepackage{tikz}
\usepackage{tasks}

\geometry{
    top=1.5cm,
    bottom=1.5cm,
    left=1.7cm,
    right=1.7cm
}

\setlength{\parindent}{0pt}

\newcommand{\version}[1]{
    \begin{center}
        \textbf{\large INTERROGATION -- Fonctions et proportions}\\
        \textbf{Version #1}
    \end{center}

    \vspace{-0.1cm}

    \textbf{Nom :} \rule{5cm}{0.4pt}
    \hfill
    \textbf{Prénom :} \rule{4cm}{0.4pt}

    \vspace{0.2cm}
}

\begin{document}

% =========================================================
% VERSION A
% =========================================================

\version{A}

\textbf{Exercice 1 -- DC \hfill /2}

Calculer $3-(-4)$

\vspace{0.1cm}

\textbf{Exercice 2 -- Tableau de valeurs \hfill /4}

On considère la fonction \(f\) définie sur l'intervalle \([0\ ;\ 4]\) par \(f(x)=x^2-3x+3\). 


\begin{multicols}{2}
    \begin{minipage}{0.45\textwidth}
    \begin{enumerate}
    \item Complétez le tableau de valeurs suivant.
    
    \begin{center}
        \renewcommand{\arraystretch}{1.5} %donne la distance entre les lignes%
        \setlength{\tabcolsep}{0.3cm} %donne la distance entre les colonnes%
        \begin{tabular}{|*{6}{c|}}
        \hline
        \(x\) & 0 & 1 & 2 & 3 & 4  \\
        \hline
        \(f(x)\) &  &  &  &  &  \\
        \hline
        \end{tabular}
    \end{center}

    \vfill\null
    
    \item À l'aide du tableau, tracer la courbe représentative de la fonction.
    \end{enumerate}
    \end{minipage}

        \hspace{-1cm}\begin{tikzpicture}[yscale=0.5, xscale=1.3]
        \draw[step=1.0,black!30!white,thin] (0,0) grid (4.5,9.5);
        \draw[->] (-0.1,0) -- (4.5,0) node[right] {\(x\)};
        \draw[->] (0,-0.5) -- (0,10) node[left] {\(f(x)\)};
        \foreach \x in {0,1,2,3,4} \draw (\x,0.5) -- (\x,-0.5) node[below] {\(\x\)};
        \foreach \y in {0,3,...,9} \draw (0.1,\y) -- (-0.1,\y) node[left] {\(\y\)};
    \end{tikzpicture}

\end{multicols}

\textbf{Exercice 3 -- Proportions \hfill /4}

Calculer :\begin{tasks}(2)
\task $50\%$ de $80$
\task $25\%$ de $75$
\end{tasks}

\vfill

% =========================================================
% VERSION B
% =========================================================

\version{B}
\textbf{Exercice 1 -- DC \hfill /2}

Calculer $-3-4$

\vspace{0.1cm}

\textbf{Exercice 2 -- Tableau de valeurs \hfill /4}

On considère la fonction \(f\) définie sur l'intervalle \([0\ ;\ 4]\) par \(f(x)=x^2-5x+7\). 


\begin{multicols}{2}
    \begin{minipage}{0.45\textwidth}
    \begin{enumerate}
    \item Complétez le tableau de valeurs suivant.
    
    \begin{center}
        \renewcommand{\arraystretch}{1.5} %donne la distance entre les lignes%
        \setlength{\tabcolsep}{0.3cm} %donne la distance entre les colonnes%
        \begin{tabular}{|*{6}{c|}}
        \hline
        \(x\) & 0 & 1 & 2 & 3 & 4  \\
        \hline
        \(f(x)\) &  &  &  &  &  \\
        \hline
        \end{tabular}
    \end{center}

    \vfill\null
    
    \item À l'aide du tableau, tracer la courbe représentative de la fonction.
    \end{enumerate}
    \end{minipage}

        \hspace{-1cm}\begin{tikzpicture}[yscale=0.5, xscale=1.3]
        \draw[step=1.0,black!30!white,thin] (0,0) grid (4.5,9.5);
        \draw[->] (-0.1,0) -- (4.5,0) node[right] {\(x\)};
        \draw[->] (0,-0.5) -- (0,10) node[left] {\(f(x)\)};
        \foreach \x in {0,1,2,3,4} \draw (\x,0.5) -- (\x,-0.5) node[below] {\(\x\)};
        \foreach \y in {0,3,...,9} \draw (0.1,\y) -- (-0.1,\y) node[left] {\(\y\)};
    \end{tikzpicture}

\end{multicols}

\textbf{Exercice 3 -- Proportions \hfill /4}

Calculer :\begin{tasks}(2)
\task $50\%$ de $120$
\task $20\%$ de $80$
\end{tasks}

\newpage

% =========================================================
% VERSION C
% =========================================================

\version{C}
\textbf{Exercice 1 -- DC \hfill /2}

Calculer $5-7$

\vspace{0.1cm}

\textbf{Exercice 2 -- Tableau de valeurs \hfill /4}

On considère la fonction \(f\) définie sur l'intervalle \([0\ ;\ 4]\) par \(f(x)=x^2-3x+4\). 


\begin{multicols}{2}
    \begin{minipage}{0.45\textwidth}
    \begin{enumerate}
    \item Complétez le tableau de valeurs suivant.
    
    \begin{center}
        \renewcommand{\arraystretch}{1.5} %donne la distance entre les lignes%
        \setlength{\tabcolsep}{0.3cm} %donne la distance entre les colonnes%
        \begin{tabular}{|*{6}{c|}}
        \hline
        \(x\) & 0 & 1 & 2 & 3 & 4  \\
        \hline
        \(f(x)\) &  &  &  &  &  \\
        \hline
        \end{tabular}
    \end{center}

    \vfill\null
    
    \item À l'aide du tableau, tracer la courbe représentative de la fonction.
    \end{enumerate}
    \end{minipage}

        \hspace{-1cm}\begin{tikzpicture}[yscale=0.5, xscale=1.3]
        \draw[step=1.0,black!30!white,thin] (0,0) grid (4.5,9.5);
        \draw[->] (-0.1,0) -- (4.5,0) node[right] {\(x\)};
        \draw[->] (0,-0.5) -- (0,10) node[left] {\(f(x)\)};
        \foreach \x in {0,1,2,3,4} \draw (\x,0.5) -- (\x,-0.5) node[below] {\(\x\)};
        \foreach \y in {0,3,...,9} \draw (0.1,\y) -- (-0.1,\y) node[left] {\(\y\)};
    \end{tikzpicture}

\end{multicols}

\textbf{Exercice 3 -- Proportions \hfill /4}

Calculer :\begin{tasks}(2)
\task $50\%$ de $400$
\task $25\%$ de $60$
\end{tasks}
\vfill

% =========================================================
% VERSION D
% =========================================================

\version{D}

\textbf{Exercice 1 -- DC \hfill /2}

Calculer $-5-3$

\vspace{0.1cm}

\textbf{Exercice 2 -- Tableau de valeurs \hfill /4}

On considère la fonction \(f\) définie sur l'intervalle \([0\ ;\ 4]\) par \(f(x)=x^2-5x+8\). 


\begin{multicols}{2}
    \begin{minipage}{0.45\textwidth}
    \begin{enumerate}
    \item Complétez le tableau de valeurs suivant.
    
    \begin{center}
        \renewcommand{\arraystretch}{1.5} %donne la distance entre les lignes%
        \setlength{\tabcolsep}{0.3cm} %donne la distance entre les colonnes%
        \begin{tabular}{|*{6}{c|}}
        \hline
        \(x\) & 0 & 1 & 2 & 3 & 4  \\
        \hline
        \(f(x)\) &  &  &  &  &  \\
        \hline
        \end{tabular}
    \end{center}

    \vfill\null
    
    \item À l'aide du tableau, tracer la courbe représentative de la fonction.
    \end{enumerate}
    \end{minipage}

        \hspace{-1cm}\begin{tikzpicture}[yscale=0.5, xscale=1.3]
        \draw[step=1.0,black!30!white,thin] (0,0) grid (4.5,9.5);
        \draw[->] (-0.1,0) -- (4.5,0) node[right] {\(x\)};
        \draw[->] (0,-0.5) -- (0,10) node[left] {\(f(x)\)};
        \foreach \x in {0,1,2,3,4} \draw (\x,0.5) -- (\x,-0.5) node[below] {\(\x\)};
        \foreach \y in {0,3,...,9} \draw (0.1,\y) -- (-0.1,\y) node[left] {\(\y\)};
    \end{tikzpicture}

\end{multicols}

\textbf{Exercice 3 -- Proportions \hfill /4}

Calculer :\begin{tasks}(2)
\task $50\%$ de $600$
\task $20\%$ de $35$
\end{tasks}

\end{document}

